\documentclass[12pt]{article}
\usepackage[utf8]{inputenc}
\usepackage[T1]{fontenc}
\usepackage{hyperref}
\title{Tasmanian, Northeastern Language Resources for K-12 Curriculum Development}
\author{Compiled by Claude Sonnet 5.5 (claude-sonnet-5-5) for Andr\'as Kornai}
\date{}
\begin{document}
\maketitle

\section{Introduction}
Northeastern Tasmanian (ISO 639-3 xpb, Glottocode benl1235) is an extinct language isolate of Tasmania. One classification treats it as a group of three varieties: Northeastern, Ben Lomond and North Midlands. Only fragmentary word lists survive, and some of them are of uncertain origin. The Tasmanian Aboriginal Centre's revival language, palawa kani, draws heavily on Northeastern and Eastern sources. My searches did not find a stand-alone Northeastern Tasmanian word list online.

\section{Primary Educational Resources}
The searches found no primers, textbooks or school programmes for Northeastern Tasmanian specifically. palawa kani is the revival effort relevant to Tasmanian schools. The Australian embassy page below describes it: \url{https://australien.ahk.de/medien/news-details/indigenous-language-series-8-palawa-kani-tasmania-1}

\section{Linguistic and Reference Materials}
\begin{itemize}
\item Tasmanian languages (Wikipedia): \url{https://en.wikipedia.org/wiki/Tasmanian_languages}
\item Eastern Tasmanian languages (Wikipedia): \url{https://en.wikipedia.org/wiki/Eastern_Tasmanian_languages}
\item Britannica, Tasmanian languages: \url{https://www.britannica.com/print/article/583948}
\end{itemize}
The main sources behind palawa kani are Plomley's 1976 word lists and Crowley and Dixon's 1981 chapter, but I did not find those texts online.

\section{Religious and Community Resources}
palawa kani is run by the Tasmanian Aboriginal Centre. The searches indicate the Centre restricts access to its language materials, so contact and permission are needed before any use. Overview: \url{https://en.wikipedia.org/wiki/Palawa_kani}

\section{Digital and Audio Resources}
No public audio, video or software for Northeastern Tasmanian was found.

\section{Curriculum Development Context}
There are no native speakers, and the surviving data are scarce and sometimes of doubtful provenance. A school unit should be built with the Tasmanian Aboriginal community, around the history of the languages and of palawa kani. It should not teach reconstructed Northeastern vocabulary on its own authority. Some critics have questioned the methods of the reconstruction, and this should be acknowledged.
\end{document}